\documentclass{beamer}
%\documentclass[17pt]{beamer}
%**********************************************************************************************************************************************
\usepackage{multicol, stmaryrd, amsfonts, graphicx, times, epsfig, amsmath, mathtools, subfigure, balance, array, siunitx, pgfgantt, xspace}
\usepackage{multirow, epsf, amsmath, amssymb, indentfirst, verbatim, keyval, url, textcomp, enumerate, calc, makecell, subfigure, xcolor, listings}
%**********************************************************************************************************************************************
\usepackage[utf8]{inputenc}
\usepackage[export]{adjustbox}
\usepackage[absolute,overlay]{textpos}
\usepackage{hyperref}
\usepackage[normalem]{ulem}
\usepackage{textpos}
\usepackage{tcolorbox}
\usepackage{listings}
%**********************************************************************************************************************************************
\useunder{\uline}{\ulined}{}%
\DeclareUrlCommand{\bulurl}{\def\UrlFont{\ttfamily\color{blue}\ulined}}
\usefonttheme{serif}
\usetheme{Pittsburgh}
\usecolortheme{beaver}
\graphicspath{ {figs/} }
\setbeamertemplate{itemize items}[ball]
\setbeamertemplate{bibliography item}[text]
\hypersetup{
   colorlinks   = true,                               %Colours links instead of ugly boxes
   urlcolor     = blue,                               %Colour for external hyper links
   linkcolor    = blue,                               %Colour of internal links
   citecolor    = red,                                %Colour of citations
   setpagesize  = false,
   linktocpage  = true,
}
%**********************************************************************************************************************************************
%\usepackage{stmaryrd, amsfonts, graphicx, times, epsfig, amsmath, mathtools,subfigure, balance, array,siunitx}
%\usefonttheme{structuresmallcapsserif}
%\usepackage{bookman}
%\usetheme{Madrid}
%\usetheme{Szeged}
%**********************************************************************************************************************************************

 %\AtBeginSection[]
%{
%  \begin{frame}
%    \frametitle{Table of Contents}
%    \tableofcontents[currentsection]
%  \end{frame}
%}

%**********************************************************************************************************************************************
%Information to be included in the title page:
\title[] %optional
{Graph Neural Network \\Lecture 3
}
%\subtitle{Graph Neural Network \\Lecture}
%**********************************************************************************************************************************************
%\author[Arthur, Doe] % (optional, for multiple authors)
%{A.~B.~Arthur\inst{1} \and J.~Doe\inst{2}}
%
%\institute[VFU] % (optional)
%{
%  \inst{1}%
%  Faculty of Physics\\
%  Very Famous University
%  \and
%  \inst{2}%
%  Faculty of Chemistry\\
%  Very Famous University
%}
%**********************************************************************************************************************************************
\date[] % (optional)

%\logo{\includegraphics[width=1in, height=0.3in]{GW_logo.eps}}
\titlegraphic{\includegraphics[width=0.8in, height=0.8in]{FrontCover_ImgID1_mod3.eps}}
%>>>>>>>>>>>>>>>>>>>>>>>>>>>>>>>>>>>>>>>>>>>>>>>>>>>>>>>>>>>>>>>>>>>>>>>>>>>>>>>>>>>>>>>>>>>>>>>>>>>>>>>>>>>>>>>>>>>>>>>>>>>>>>>>>>>>>>>>>>>>>>>
%>>>>>>>>>>>>>>>>>>>>>>>>>>>>>>>>>>>>>>>>>>>>>>>>>>>>>>>>>>>>>>>>>>>>>>>>>>>>>>>>>>>>>>>>>>>>>>>>>>>>>>>>>>>>>>>>>>>>>>>>>>>>>>>>>>>>>>>>>>>>>>>


\begin{document}

\begin{frame}
    % Print the title page as the first slide
    \titlepage
\end{frame}

\begin{frame}{Overview}
    % Throughout your presentation, if you choose to use \section{} and \subsection{} commands, these will automatically be printed on this slide as an overview of your presentation
    \tableofcontents
\end{frame}

%------------------------------------------------
\section{Introduction to GCNs}

%------------------------------------------------
\begin{frame}{What is a Graph Convolutional Network (GCN)?}
    \begin{itemize}
        \item A type of Graph Neural Network (GNN) designed for graph-structured data.
        \item Extends the concept of convolutional neural networks (CNNs) to graphs.
        \item Captures local graph structure and node features.
    \end{itemize}
\end{frame}


%------------------------------------------------
\begin{frame}{Why GCNs?}
    \begin{itemize}
        \item \textbf{Graph-Structured Data}: Many real-world problems involve graphs (e.g., social networks, molecular structures).
        \item \textbf{Node-Level Tasks}: Node classification, link prediction.
        \item \textbf{Graph-Level Tasks}: Graph classification, clustering.
    \end{itemize}
\end{frame}

%------------------------------------------------
\section{Mathematical Formulation}
\begin{frame}{GCN Layer}
    The GCN layer is defined as:
    \[
    H^{(l+1)} = \sigma\left(\tilde{D}^{-\frac{1}{2}} \tilde{A} \tilde{D}^{-\frac{1}{2}} H^{(l)} W^{(l)}\right)
    \]
    \begin{itemize}
        \item \(\tilde{A} = A + I\): Adjacency matrix with self-loops.
        \item \(\tilde{D}\): Degree matrix of \(\tilde{A}\).
        \item \(H^{(l)}\): Node features at layer \(l\).
        \item \(W^{(l)}\): Learnable weight matrix.
        \item \(\sigma\): Activation function (e.g., ReLU).
    \end{itemize}
\end{frame}

%------------------------------------------------
\begin{frame}{Intuition Behind GCN}
    \begin{itemize}
        \item \textbf{Normalization}: \(\tilde{D}^{-\frac{1}{2}} \tilde{A} \tilde{D}^{-\frac{1}{2}}\) ensures normalized aggregation.
        \item \textbf{Feature Propagation}: Each node aggregates features from its neighbors.
        \item \textbf{Nonlinearity}: Activation function introduces nonlinearity.
    \end{itemize}
\end{frame}
%------------------------------------------------




%----------------------------------------------------------------------------------------

\end{document}